% 换测度前后的联合分布与乘积分布。
% 设定：z1 等概率取 0、1；算法以概率 1-q 输出 W=z1、以概率 q 输出 1-z1，q=0.1。
% 联合分布 P(z1,W)：对角 0.45，非对角 0.05；乘积分布 P(z1)P(W)：每格 0.25。
% 互信息 I(W;z1)=sum P log(P/Q)=2*0.45*log(1.8)+2*0.05*log(0.2)=log2-h(0.1)≈0.368 nat。
% 右图：连续情形 n=1，z1~U[0,1]，精确输出样本均值 W=z1。
% 联合分布全部质量在对角线 w=z1 上（面积为零），乘积分布为正方形上的均匀分布，
% 后者给对角线的概率为 0，故 KL 与互信息为无穷。此面板只作示意。
\begingroup
\providecolor{FigureGreen}{HTML}{829C8B}
\providecolor{FigureGreenFill}{HTML}{E8EEE9}
\begin{tikzpicture}[figure picture,x=1mm,y=1mm]
\def\cell#1#2#3#4{% x y value fill-shade
 \fill[FigureBlue!#4!white] (#1,#2) rectangle ({#1+12},{#2+12});
 \draw[FigureStroke,line width=.5pt] (#1,#2) rectangle ({#1+12},{#2+12});
 \node at ({#1+6},{#2+6}) {$#3$};}
% 联合分布
\node[figure heading] at (24,38) {联合分布 $P_{z_1,W}$};
\cell{12}{12}{0.45}{55}\cell{24}{12}{0.05}{6}
\cell{12}{0}{0.05}{6}\cell{24}{0}{0.45}{55}
% 坐标: 上行 z1=0，下行 z1=1；左列 W=0，右列 W=1
\node[figure note] at (18,26.5) {$W=0$};\node[figure note] at (30,26.5) {$W=1$};
\node[figure note,anchor=east] at (11,18) {$z_1=0$};\node[figure note,anchor=east] at (11,6) {$z_1=1$};
% 乘积分布
\node[figure heading] at (68,38) {乘积分布 $P_{z_1}\otimes P_W$};
\cell{56}{12}{0.25}{28}\cell{68}{12}{0.25}{28}
\cell{56}{0}{0.25}{28}\cell{68}{0}{0.25}{28}
\node[figure note] at (62,26.5) {$W=0$};\node[figure note] at (74,26.5) {$W=1$};
\draw[FigureStroke,line width=.8pt,-{Stealth[length=1.6mm]}] (40,12)--(52,12);
\node[figure note] at (46,15.5) {消除依赖};
\node[figure note] at (46,-5.5) {$I(W;z_1)=\mathrm{KL}(P_{z_1,W}\,\|\,P_{z_1}\otimes P_W)\approx0.368$};
% 连续情形示意
\node[figure heading] at (113,38) {连续输出 $W=z_1$};
\fill[FigureBlueFill] (101,0) rectangle (125,24);
\draw[FigureStroke,line width=.5pt] (101,0) rectangle (125,24);
\draw[FigureBlue,line width=1.6pt] (101,0)--(125,24);
\node[figure note] at (113,26.5) {横轴 $z_1$，纵轴 $W$};
\node[figure note] at (113,-5.5) {浅底：乘积分布；对角线：联合分布};
\end{tikzpicture}
\endgroup
